\section{Burst time and burst size}
\label{sec:burst-time-size}

The two random quantities that a bursting cell presents to the outside world
are the time of its rupture and the size of the rupture. Both are now within
reach of closed forms.

\subsection{Burst time}
\label{sec:burst-time}

The no-burst probability is $I(t)=\Prob{\tau>t}$, so the density of the
burst time is
\begin{equation}
\varphi(t):=-I'(t)=\delta J(t),
\label{eq:phi}
\end{equation}
using $I'=-\delta J$ from \S\ref{sec:moments-recap}. This density is
\emph{defective}: since $I(t)\to b$,
\begin{equation}
\int_0^\infty\varphi(t)\,\dd t=1-b,
\end{equation}
with the remaining mass $b$ sitting at $\tau=\infty$ --- the realisations
that clear internally without ever bursting. Conditioned on eventual burst,
the density is $\varphi(t)/(1-b)=\delta J(t)/(1-b)$. Where earlier work
referred to ``the known pdf of rupture times $\phi(t)$'', it is
\eqref{eq:phi} that was meant, available all along from $I'=-\delta J$.

\subsection{Burst size}
\label{sec:burst-size}

The burst size $\Kb$ is the value of $\xt$ immediately before rupture. For
the burst to occur in state $k$, the process must be in state $k$ and then
rupture from there; rupturing from state $k$ happens at rate $\delta k$.
Hence
\begin{equation}
\pr{\Kb = k} = \delta k \int_0^\infty p_k(t)\,\dd t.
\label{eq:burstintegral}
\end{equation}
A direct discretisation makes this plain: partitioning time into intervals
of width $\Delta t$, the probability that the burst occurs during an interval
in which the process is at $k$ is
$$
\pr{\xt=k}\,\pr{\text{rupture in }\Delta t\mid\xt=k}
=\pr{\xt=k}\,\delta k\,\Delta t,
$$
and summing over intervals and passing to the limit $\Delta t\to0$ yields
\eqref{eq:burstintegral}.

The integral is evaluable in closed form. With $P(t)$ as in \eqref{eq:Pdef},
\begin{equation}
\frac{\dd}{\dd t}\Bigl[\frac{P(t)^k}{k\beta}\Bigr]
=\frac{1}{\beta}P(t)^{k-1}P'(t)
=\lrb{\frac{a-b}{b-a\ee^{\beta(a-b)t}}}^2
\ee^{\beta(a-b)t}P(t)^{k-1}
=p_k(t),
\label{eq:antideriv}
\end{equation}
as one verifies from
$P'(t)=\beta(a-b)^2\ee^{\beta(a-b)t}/(b-a\ee^{\beta(a-b)t})^2$.
Therefore, since $P(0)=0$ and $P(\infty)=1/a$,
\begin{equation}
\int_0^\infty p_k(t)\,\dd t=\frac{1}{k\beta}\Bigl(\frac{1}{a}\Bigr)^k,
\end{equation}
and \eqref{eq:burstintegral} gives the burst-size law.

\begin{theorem}[Burst-size distribution]
\label{thm:burst}
Unconditionally (including the mass $b$ of realisations that never burst),
\begin{equation}
\pr{\Kb=k}=\frac{\delta}{\beta}\,a^{-k},\qquad k\ge1;
\label{eq:burstlaw}
\end{equation}
conditioned on the burst occurring,
\begin{equation}
\pr{\Kb=k\mid\text{burst}}=(a-1)\,a^{-k},\qquad k\ge1.
\label{eq:burstlaw-cond}
\end{equation}
\end{theorem}

\begin{proof}
The unconditional form is the computation above. For the conditional form,
divide \eqref{eq:burstlaw} by the eventual-burst probability $1-b$:
$$
\pr{\Kb=k\mid\text{burst}}
=\frac{(\delta/\beta)\,a^{-k}}{1-b}
=\frac{AB}{B}\,a^{-k}
=(a-1)\,a^{-k},
$$
using $\delta/\beta=AB$ from \eqref{eq:AB}.
\end{proof}

No simulation is required; nevertheless two consistency checks close the
loop pleasantly. First, normalisation:
$$
\sum_{k=1}^{\infty}\frac{\delta}{\beta}a^{-k}
=\frac{\delta}{\beta}\,\frac{1}{a-1}
=\frac{AB}{A}=B=1-b,
$$
matching the eventual-burst probability. Second, the first moment:
$$
\sum_{k=1}^{\infty}k\,\frac{\delta}{\beta}a^{-k}
=\frac{\delta}{\beta}\,\frac{a}{(a-1)^2}
=\frac{aB}{A}=V_\infty,
$$
recovering \eqref{eq:Vinf}. This second check independently validates the
burst-size law, the formula for $V_\infty$, and the identity
\eqref{eq:AB} simultaneously.

\subsection{Burst size equals the quasi-stationary distribution}
\label{sec:burst=qsd}

Compare \eqref{eq:burstlaw-cond} with Theorem~\ref{thm:qsd}: they are the
same law.

\begin{corollary}[Burst size $=$ QSD]
\label{cor:burst-qsd}
The burst-size distribution, conditioned on bursting, is exactly the
quasi-stationary distribution of the load: both are geometric on
$\{1,2,\ldots\}$ with ratio $1/a$.
\end{corollary}

This is not obvious, and it deserves a moment's pause. A burst is a
\emph{size-biased} draw from the transient load distribution, taken against
the burst intensity, and then integrated over all burst times; that this
should land exactly on the conditional limiting distribution is a genuine
identity, not a definition. The algebra above proves it; a structural
argument --- presumably relating size-biasing composed with the burst
intensity to the Yaglom limit \cite{yaglom1947certain} --- would make it
quotable rather than coincidental, and remains open
(see \S\ref{sec:open-problems}).

\subsection{Size-biasing and the size of a late burst}
\label{sec:size-bias}

The size-biasing mechanism is worth stating on its own terms. A cell bursts
at rate $\delta\xt$, so given that the burst occurs at time $t$, its size is
a size-biased draw from $p_\cdot(t)$:
\begin{equation}
\pr{\Kb=k\mid\tau=t}=\frac{k\,p_k(t)}{J(t)},
\qquad
\mean{\Kb\mid\tau=t}=\frac{K(t)}{J(t)}
=1+\frac{2\beta}{\delta}\lrb{1-I(t)}.
\label{eq:sizebias}
\end{equation}
Taking $t\to\infty$,
\begin{equation}
\mean{\Kb\mid\tau=t}\xrightarrow[t\to\infty]{}
1+\frac{2\beta}{\delta}(1-b)
=\frac{\langle X^2\rangle_{\rm QS}}{\langle X\rangle_{\rm QS}},
\end{equation}
where the second equality follows from \eqref{eq:QSmoments}. A very late
burst is roughly twice the size of an average one --- the signature of
size-biasing a geometric law: the cells that survive long are precisely the
ones that grew large, and large cells burst quickly.
